% page 274 du fac-simile (image p-273 du scan)
% bloc 154.9 x 242.0 mm ; marge gauche 8.0 mm
\pgc{-12.700mm}{0.000mm}{
\l{\cel{3}266}
\l{}
\l{\sou{karezar}.\cel{2}(\sou{trans.}) Tushar lejere per sua manuo, sua labii,}
\l{\cel{3}por signifikar sua afeciono, tenereso, amo, amoro - ed,}
\l{\cel{3}en ula kazi, por produktar efekti erotika. - EFIS.}
\l{}
\l{\sou{kargar}. (\sou{trans}.)\cel{2}Igar (ulu, ulo) subisar la pezo da objek-}
\l{\cel{3}to qua esas portenda ad altra loko, transportenda. EFIS.}
\l{}
\l{\sou{kariar}. (\sou{netrans}.)\cel{2}(\sou{patol}.) I. (\sou{aludante osto}). Domajesar}
\l{\cel{3}(da morbo) forme di molesko di la tisuo osta kun pusifo.}
\l{\cel{3}II. (\sou{aludante la denti}). Domajesar (da morbo) forme\cel{2}di}
\l{\cel{3}molesko lenta e gradoza, pri emalio ed ivoro. - DEFIS.}
\l{}
\l{\sou{kariatido}.\cel{2}Statuo de muliero, drapirita, quan onu reprezen-}
\l{\cel{3}tas kom stacanta, e qua subtenas entablamento,kornico,}
\l{\cel{3}vice kolono, pilastro. - DEFIS.}
\l{}
\l{\sou{karibuo}. (\sou{zool.)}\cel{2}Rentiro di Kanada. L.\cel{1}\sou{Rangifer caribus}.}
\l{\cel{3}- DEF.}
\l{}
\l{\sou{karico}. (\sou{bot.)}\cel{2}Planto ek la familio \textquotedbl{}ciperaci\textquotedbl{}. L.\cel{1}\sou{carex}.}
\l{\cel{3}EFIS.}
\l{}
\l{\sou{kariero}. Profesiono ofico-hierarkioza. - DEFIRS.}
\l{}
\l{\sou{kariko}.\cel{2}Redingoto ampla, kun pelerino dispozita etajatre.DF}
\l{}
\l{\sou{karikatar}. (\sou{trans.)}\cel{3}Reprezentar persono, kozi, ma exajer-}
\l{\cel{3}ante, groteske, por igar li rido-mokinda. - DEFIS.}
\l{}
\l{\sou{kariocinezo}. (\sou{biol.}) Divido di la celulo, lor qua la divido}
\l{\cel{3}di la nuklei preiras olta di la celulo-korpo.}
\l{}
\l{\sou{kariofilo}.\cel{2}(\sou{bot.}) Burjono di la floro di arboro.(L.\cel{1}\sou{caryo}-}
\l{\cel{3}\sou{phillus aromaticus}) di Moluko o diAntili,ek la familio}
\l{\cel{3}\textquotedbl{}mirtacei\textquotedbl{},uzata kom spico, piklo. -\cel{2}I.}
\l{}
\l{\sou{kariolizo}. (\sou{bot.})\cel{2}Dissolvo di la celulo-nukleo.}
\l{}
\l{\sou{karitato}. I. Amo kompatoza di la kunhomi. - II. (\sou{kristanismo})}
\l{\cel{3}Amo di la kunhomo, por deo. - EFIS.}
\l{}
\l{\sou{karlino}. (\sou{bot.}) Planto di qua la radiko uzesis kom sudorif-}
\l{\cel{3}igivo\cel{2}- L.\cel{1}\sou{carlina}. - EFISL.}
\l{}
\l{\sou{karmelito}.\cel{2}(\sou{religio katol.}) Reguliero ek la ordeno \textquotedbl{}karmel\textquotedbl{}}
\l{\cel{3}(ordeno fondita en 1451,da Jean Soreth).}
\l{}
\l{\sou{karmezino}.\cel{2}Koloro obskure-reda briloze. - DEFIS.}
\l{}
\l{\sou{karmino}. Koloro bele-reda, ma min brilo-reda e plu brune-reda}
\l{\cel{3}kam cinabrea. - DEFIRS.}
\l{}
\l{\sou{karno.}\cel{2}(\sou{anat.}) La muskuli di la korpo di la homo o di anima-}
\l{\cel{3}lo (kontraste kun la osti e la pelo). - DEFIS.}
}
